\begin{textsample}{NEG Dim 9 – global\_north\_2024\_05 – Score -2.00 – global\_north\_2024\_05/108658655.txt}  \label{ex:f9_neg_165}
The San Jose State student encampment opposing the war in Gaza, which had been in place for over a week as the school approached its commencement, came down Wednesday evening after a meeting between the protesters and school administrators.
The coalition of students leading the encampment -- SJSU People ' s University for Gaza and Students for Justice in Palestine -- took down their tents and canopies on Wednesday night around 7:30 p.m.
It was the most recent encampment at a Bay Area university to be disbanded after such camps began popping up last month.
On Tuesday, the student encampment at the University of San Francisco also came down.
However, students at Stanford University were continuing to camp Thursday in protest of the war in Gaza.
Students broke down the camp at San Jose State following a Monday afternoon meeting between representatives from Students for Justice in Palestine and some of the university ' s administration, which included President Cynthia Teniente-Matson and Interim Vice President for Student Affairs Mari Fuentes-Martin, said university spokesperson Michelle Smith McDonald.
The camp was in place since May 13. @ @ @ @ @ @ @ @ @ @ administration, which included an acknowledgement of genocide of Palestinians in Gaza, transparency from the university on their financial investments and holdings, and divestment from companies that support Israel.
San Jose State officials responded in an email to students at the encampment, stating that the university had publicly posted a lengthy set of disclosures and document regarding investments, including by nonprofit organizations linked to the school.
The university also stated that they follow the California State University system ' s investment policy, which states that the CSU has"no direct \textbf{ownership} of stocks or bonds based in Israel."They also stated that \$3.2 million, or less than 0.04 \% of CSU investments as of March 31, were invested in mutual funds in Israel.
Teniente-Matson also suggested the creation of a student advisory council made up of students from organizations that are affiliated with Middle Eastern issues, cultures and nationalities to work with faculty and address students ' concerns of university partnerships with Israel for study-abroad programs and internships with companies affiliated with Israel. @ @ @ @ @ @ @ @ @ @ although they appreciated the progress made toward addressing their demands, they still asked the university to fully address all of their demands."This is not the end,"the students wrote in the post."While there has been progress made towards most of our demands, and commitments to continue working on them, we are disappointed with SJSU ' s overall response to this genocide these past eight months."Roughly a day before San Jose State students took down their encampment, students at the University of San Francisco abandoned their 50-tent camp, allowing university staff to pack up the tents and abandoned belongings.
The students that organized the encampment on Welch Field, which was in place since April 29, demanded that USF acknowledge the genocide of Palestinians, disclose all financial holdings the university has and divest from any companies that may be associated with Israel, end all academic partnerships with Israel and protect pro-Palestine speech and students on campus, according to a statement by university President Paul J.
Fitzgerald on May @ @ @ @ @ @ @ @ @ @ to calling for a ceasefire in Gaza.
USF also agreed to establish a Socially Responsible Investment Task Force, which will include at least one representative from the encampment.
Fitzgerald also stated that the university does not have any academic connections with Israeli institutions, but is investigating any academic opportunities and trips that may exist.
The encampment ended Monday as USF staff packed unoccupied tents.
Students who were in the process of decamping stated they were planning on leaving the same day.
In contrast, Stanford University students continued to camp Thursday and demand the university divest from companies and companies supporting Israel ' s military offensive in Gaza and call for a ceasefire.
Mark Cu, a reporter with the student newspaper Stanford Daily, said that the students are currently pressuring the university to drop the academic sanctions against students involved in the protest.
On May 13, the encampment at San Francisco State came to an end after 15 days, following an agreement between the protesting students and President Lynn Mahoney.
The university committed @ @ @ @ @ @ @ @ @ @ that will develop new screenings for investments that violate human rights and divesting from those investments,"according to a statement by Students for Gaza.
On May 14, UC Berkeley students and staff \textbf{quietly} packed up the dozens of tents that had previously stood at Sproul Plaza for months.
The students and faculty involved called for the university to disavow the war in Gaza and join calls for a permanent and immediate ceasefire, divest from all financial holdings that support Israel ' s military operations, end academic partnerships with Israeli universities and provide more support to Palestinian students and establish a Palestinian studies program.

% matched lemmas: ownership, quietly
\end{textsample}
